\chapter{Implementación}
\label{chap:implementacion}

\lettrine{E}{n} este capítulo se describen los aspectos prácticos de la implementación del sistema, con el objetivo de explicar qué software se necesita para compilarlo, cómo está organizado el código fuente y qué pasos hay que seguir para ejecutarlo en desarrollo.

\section{Software requerido}

Para compilar y ejecutar el proyecto en un entorno de desarrollo es necesario instalar manualmente el siguiente software. El resto de dependencias (Spring Boot, React, MapLibre, Firebase Client, etc.) se descargan automáticamente a través de Maven, npm o Gradle al construir cada subsistema.

\begin{table}[H]
  \centering
  \rowcolors{2}{white}{udcgray!15}
  \begin{tabular}{p{0.28\textwidth}|p{0.28\textwidth}|p{0.28\textwidth}}
    \rowcolor{udcpink!25}
    \textbf{Backend} & \textbf{Frontend} & \textbf{Android} \\
    Java 21 (Eclipse Temurin) & Node.js 18+ y npm & Android Studio \\
    Apache Maven 3.6+ & Git & Android SDK 34+ \\
    Docker y Docker Compose & -- & Git \\
    Git & -- & Configuración local de Firebase \\
  \end{tabular}
  \caption{Software mínimo necesario para levantar cada subsistema}
  \label{tab:software-requerido}
\end{table}

El backend necesita Java 21 con Maven para compilar y ejecutar el servicio. La base de datos PostgreSQL con PostGIS se levanta mediante Docker con el fichero \textit{docker-compose.yaml} incluido en el proyecto, por lo que no es necesario instalar PostgreSQL manualmente.

El frontend solo requiere Node.js 18 o superior y npm para instalar las dependencias de React, Vite y las librerías asociadas.

La aplicación Android necesita Android Studio con el SDK 34, y hay que añadir manualmente el fichero \textit{google-services.json} de Firebase para que las notificaciones push funcionen correctamente.

En los tres casos es necesario disponer de Git para clonar el repositorio. El resto de dependencias se gestionan automáticamente al compilar cada módulo, por lo que no es necesario instalarlas manualmente.

Para la puesta en producción, se recomienda configurar un entorno de despliegue con Docker para el backend y la base de datos, y utilizar servicios como Firebase Hosting para el frontend y Google Play Store para la aplicación Android. Esto viene explicado con más detalle en el apéndice \ref{sec:instalacion}, donde se incluyen las instrucciones de instalación y despliegue de cada subsistema.

\section{Estructura del repositorio}

El repositorio se organiza en tres directorios raíz, cada uno correspondiente a un subsistema: \textit{backend/}, \textit{frontend/} y \textit{EmergencyApp/}. A continuación se describe la organización dentro de cada uno de ellos y los ficheros más relevantes.

\subsection{Backend}

El directorio \textit{backend/} es un proyecto Maven convencional. El código fuente principal se encuentra en \textit{src/main/java}, los recursos (configuración de Spring, scripts DDL, plantillas) en \textit{src/main/resources}, y los tests en \textit{src/test/java}. Dentro de \textit{resources} destacan dos subdirectorios propios del proyecto: \textit{openApi/}, que contiene el contrato YAML de la API, y \textit{db/}, con los scripts de creación e inicialización de la base de datos. El fichero \textit{docker-compose.yaml} se encuentra también en \textit{resources} y permite levantar PostgreSQL con PostGIS en un contenedor. El \textit{pom.xml} declarado en la raíz gestiona las dependencias, el plugin de OpenAPI y el resto de librerías del backend.

\subsection{Frontend}

El directorio \textit{frontend/} contiene la aplicación React con Vite. El código fuente se ubica en \textit{src/}, con subdirectorios por funcionalidad: \textit{api/} para el acceso al backend mediante RTK Query, \textit{features/} para las pantallas agrupadas por dominio, \textit{app/} para la configuración del almacén Redux y utilidades compartidas, \textit{locales/} para los ficheros de traducción (ES/GL) y \textit{features/map/} para la lógica cartográfica con MapLibre. En la raíz se encuentran \textit{package.json} con las dependencias, \textit{vite.config.js} para la configuración del empaquetador y \textit{index.html} como punto de entrada.

\subsection{Aplicación Android}

El directorio \textit{EmergencyApp/} es un proyecto Gradle. El módulo principal está en \textit{app/}, que contiene \textit{src/main/java} con el código Kotlin, \textit{src/main/res} con los recursos (cadenas, iconos) y el \textit{AndroidManifest.xml}. Dentro del código fuente, los paquetes se organizan por responsabilidad técnica: \textit{ui/} para las pantallas con Jetpack Compose, \textit{net/} para el cliente HTTP con Ktor, \textit{messaging/} para la integración con Firebase Cloud Messaging, \textit{data/dto/} para los objetos de intercambio con la API y \textit{util/} para funciones auxiliares. Los ficheros \textit{build.gradle.kts} de la raíz y del módulo \textit{app/} gestionan las dependencias del proyecto.

\section{Instrucciones de compilación}

\subsection{Backend}

Para compilar y ejecutar el backend, primero hay que levantar la base de datos local y después generar y compilar el proyecto. El orden recomendado es el siguiente:

\begin{lstlisting}[language=Bash]
cd backend
docker compose -f src/main/resources/docker-compose.yaml up -d
mvn -DskipTests generate-sources compile
mvn spring-boot:run
\end{lstlisting}

El primer comando inicia PostgreSQL con PostGIS. El segundo ejecuta la generación de OpenAPI y compila el proyecto, incluyendo las fuentes generadas. El tercero arranca la aplicación en modo desarrollo. Si se quiere ejecutar la batería de pruebas backend, el comando habitual es `mvn test`.

\subsection{Frontend}

Para ejecutar la interfaz web, basta con instalar las dependencias y lanzar Vite desde el directorio \textit{frontend/}:

\begin{lstlisting}[language=Bash]
cd frontend
npm install
npm run dev
\end{lstlisting}

Este flujo inicia el servidor de desarrollo y permite trabajar sobre la aplicación sin necesidad de generar primero una compilación final. Cuando se quiera validar el resultado de producción, puede utilizarse `npm run build`.

\subsection{Aplicación Android}

La aplicación Android se abre desde Android Studio a partir del directorio \textit{EmergencyApp/}. Antes de ejecutarla es necesario configurar localmente Firebase y disponer del fichero \textit{google-services.json} correspondiente. Una vez importado el proyecto, la aplicación puede lanzarse desde el propio IDE o mediante Gradle:

\begin{lstlisting}[language=Bash]
cd EmergencyApp
.\gradlew.bat assembleDebug
\end{lstlisting}

En la práctica, la ejecución diaria se realiza desde Android Studio, ya que facilita la selección del dispositivo, la depuración y la gestión de permisos. El comando anterior resulta útil para verificar que el módulo compila correctamente desde línea de comandos.

En conjunto, esta organización permite compilar cada subsistema de forma independiente y reproducible, manteniendo separados el backend, la interfaz web y la aplicación móvil.
